% Target: rmp-workbench-ii-positive-mpo-old
% Formula: O = X^-1 (sum_k A^k tensor conjugate(A^k)) X.
% Ink: Canonical public kernel frames; no raw TikZ.
% Boundary: O has four open ports; each capsule branches one outer virtual leg into
%   the two A rails, whose physical legs form one open vertical line.
% Source: author source ImagesReview Section II/ImagesReview Section II.tex lines
%   317-335; archival output sec2fig6old.pdf
% Capabilities: kernel, typed-ports
\[
  \tenkzkernel
  \begin{tenkz}[rows={wire}, cols=1, bonds=none]
    \tn[at=(1,1), skin=box,
      ports={90:physical,270:physical,0:virtual,180:virtual}]{O}
  \end{tenkz}
  \;=\;
  \begin{tenkz}[rows={wire,wire,wire}, cols=5, bonds=none]
    \tn[at=(2,1), skin=box, name=xi,
      ports={180:virtual,45:virtual,315:virtual}]{X^{-1}}
    \tn[at=(1,3), skin=box, name=A,
      ports={180:virtual,0:virtual,90:physical,270:physical}]{A}
    \tn[at=(3,3), skin=box, name=Ab,
      ports={180:virtual,0:virtual,90:physical,270:physical}]{\overline{A}}
    % One size class gives X and X^{-1} the same measured capsule.
    \tn[at=(2,5), skin=box, name=x,
      ports={135:virtual,225:virtual,0:virtual}]{X}
    \tnwire{open w}{xi.180}
    \tnwire[route=arc]{xi.45}{A.180}
    \tnwire[route=arc]{xi.315}{Ab.180}
    \tnwire[route=arc]{A.0}{x.135}
    \tnwire[route=arc]{Ab.0}{x.225}
    \tnwire{x.0}{open e}
    \tnwire{open n}{A.90} \tnwire{A.270}{Ab.90}
    \tnwire{Ab.270}{open s}
  \end{tenkz}
\]
